\section{Evaluation design}
\label{sec:editorial-evaluation}
Appendix~\ref{app:comparison-inventory} indexes the benchmark-specific metrics,
complete method scorecards and outstanding external comparisons.
Appendix~\ref{app:metric-completion} records the remaining current-model control,
decoded-video and efficiency measurements and their prerequisites.
\paragraph{Data and matched information access.}
The selected DROID collection contains 1,126 recordings, split by recording
session into 851 training, 143 validation and 132 original test episodes.
The spatial study uses only training and validation. At ten query steps,
141 validation episodes from 59 sessions yield 1,631 eligible windows.
All DROID scores use the common training-standardized $4\times4$ grid;
pooled $2\times2$ coordinates are supplementary. Compared predictors share
features, history and command access except for each explicitly removed path.
Persistence copies the last observation.

The principal reported endpoint is absolute native-grid standardized feature MSE
at query step ten. For a comparator with error $e_b$ and ShiftWM error $e_o$,
we additionally report $100(e_b-e_o)/e_b$ as relative error reduction; this is not a
percentage-point change in task success. Five-step endpoints and the full
forecast curves remain visible as secondary diagnostics.

\paragraph{Training and uncertainty.}
The initial campaign trains five arms with three seeds for 30 epochs.
A registered follow-up adds six models to complete the mixing-by-bounding
comparison after the initial controls were revealed. All 21 runs complete
the full budget. Checkpoints minimize window-weighted validation loss across
all ten query steps. Reported endpoint scores average windows within episodes,
then weight episodes and seeds equally. Paired intervals resample sessions
and training seeds with 10,000 draws. These are exploratory development
comparisons without multiplicity correction; the ten-step endpoint is not
an untouched confirmatory test. All horizons, 36 paired-bootstrap contrasts
(32 method comparisons and four interaction effects), selected checkpoints
and exact reload checks appear in Appendix~\ref{app:spatial-development}.

\paragraph{Single-observation transfer.}
The IWS study trains separate task adapters on recorded PushT, Box and Rope
trajectories. Each receives one image and 60 native commands and predicts
the next 59 feature grids, with no observed-transition context or FiLM.
Three learned arms, three tasks and three seeds give 27 full-budget runs;
persistence is untrained. A component follow-up adds nine no-$\tanh$ models.
Checkpoints minimize development endpoint MSE; all 36 were fixed before reserved
payload access. The reserved population has 600 original handles from 30
trajectories, with no checkpoint reselection from its outcomes.
The primary comparison is bounded mixing versus additive anchoring, with
paired seed--trajectory intervals and an equal-task mean of relative gains.
Appendix~\ref{app:experiment-alignment} gives the complete input, split and
selection contracts and the disclosed numerical execution revision. Development
and reserved upstream-validation outcomes use distinct complete-study finalizers.

\paragraph{Forecast metrics and task outcomes.}
DROID and IWS report standardized MSE, standardized MAE, raw DINOv2-feature L1 and cosine distance. The DROID additions use all fixed checkpoints and the original development windows (Table~\ref{tab:droid-complementary-metrics}); its primary endpoint remains unchanged.

These absolute errors and their relative
reductions assess recorded-action forecasting, not closed-loop success.
RLA-WM \citep{zhang2026rlawm} also evaluates DINO-token L1 and decoder-based
LPIPS/SSIM; the latter require RGB predictions, which our model does not produce.
Closed-loop planning instead requires task-outcome evaluation, as in LeWM and
DINO-WM \citep{maes2026lewm,zhou2024dinowm}.

\paragraph{Historical studies remain distinct.}
The earlier two-context model has different supervision and representations.
Its simulation, original real-video test and fresh-session confirmation are
preserved in Appendices~\ref{app:original-context-study},
\ref{app:real-video} and~\ref{app:fresh-real-video}. Their primary outcomes and
mixed findings are unchanged. They are neither ablations of this spatial
predictor nor evidence of its cross-domain or closed-loop performance.
